%% Based on aiptemplate.tex for RevTeX 4.1 / AIP substyles
%%
%% Gauge coupling constants from tetrahedral combinatorics
%% B. W. Mayes — April 2026
%%

\documentclass[aip,graphicx]{revtex4-1}

\usepackage{amsmath,amssymb}

\begin{document}

\title{Gauge coupling constants from tetrahedral combinatorics}

\author{Benjamin Walker Mayes}
\email[Corresponding author: ]{bmayes@unified-dynamics.org}
\affiliation{Unaffiliated, Cedar Park, Texas 78613, USA}

\author{Arthur Eugene Mayes}
\affiliation{Loyola University School of Law, USA}

\author{Dustin Huff}
\affiliation{Unaffiliated, USA}

\author{Brayden Ross Sanders}
\affiliation{Unaffiliated, USA}

\author{Jonathan Engle}
\affiliation{Unaffiliated, USA}

\author{Jeffrey A. Dewey}
\affiliation{Unaffiliated, USA}

\author{Brandon Bevin}
\affiliation{Georgia Southern University, USA}

\author{A. Christopher Beckingham}
\affiliation{Queen's University, Canada}

\date{\today}

\begin{abstract}
We derive the strong coupling constant $\alpha_s(m_Z)$, the weak
mixing angle $\sin^2\theta_W(m_Z)$, and the inverse fine structure
constant $\alpha_{\mathrm{EM}}^{-1}(m_Z)$ from the combinatorial
geometry of the complete graph $K_4$ with zero free parameters.
The Standard Model gauge group
$\mathrm{SU}(3)\times\mathrm{SU}(2)\times\mathrm{U}(1)$ emerges
from the face, edge, and vertex structure of $K_4$. Coupling ratios
follow from boundary element counting; the overall normalization is
fixed by the Regge deficit angle on the $K_4$ two-skeleton; and the
matter sector coupling is suppressed relative to the geometry coupling
by $1/d^3 = 1/27$, where $d = V-1 = 3$ is the vertex degree.
We further show that the three charged-lepton generations correspond
to the three cosets of the Klein four-group $V_4$ in the alternating
group $A_4 \cong A_4/V_4 \times V_4$, with $A_4/V_4 \cong Z_3$, and
that the Koide mass formula is the $Z_3$ character expansion of the
lepton mass spectrum---a consequence of group structure, not a
numerical coincidence.
We obtain $\alpha_s(m_Z) = 129/(112\pi^2) = 0.1167$ (measured
$0.1180 \pm 0.0009$; deviation 1.1\%),
$\sin^2\theta_W(m_Z) = 3/13 = 0.2308$ (measured $0.23121
\pm 0.00004$; deviation 0.19\%), and
$\alpha_{\mathrm{EM}}^{-1}(m_Z) = 13\pi^2 = 128.3$ (measured
$127.951 \pm 0.009$; deviation 0.28\%).
\end{abstract}

\pacs{12.10.-g, 11.15.Ha, 02.40.Re, 02.10.Ox}

\maketitle


% ============================================================
\section{Introduction}
\label{sec:intro}
% ============================================================

The Standard Model of particle physics contains 19 free parameters
whose values are determined by experiment rather than by the theory
itself.\cite{pdg2024} Among these, the gauge coupling constants
$\alpha_s$, $\sin^2\theta_W$, and $\alpha_{\mathrm{EM}}$ set the
strengths of the strong, weak, and electromagnetic interactions.
Despite decades of effort in grand unification
\cite{georgi1974,pati1974} and string theory,\cite{polchinski1998}
no framework has derived these constants from geometry alone with
zero adjustable parameters.

We present such a derivation. The complete graph $K_4$---four
vertices connected by all six edges, forming a tetrahedron---serves
as the sole geometric input. The Standard Model gauge group, coupling
ratios, and overall normalization all follow from the combinatorial
properties of $K_4$ and the Regge calculus on its two-skeleton. No
renormalization group running, matching scale, or fitted parameter
enters the calculation.


% ============================================================
\section{The $K_4$ framework}
\label{sec:framework}
% ============================================================

The complete graph $K_4$ has $V = 4$ vertices, $E = \binom{4}{2}= 6$
edges, $F = \binom{4}{3} = 4$ triangular faces, and Euler
characteristic $\chi = V - E + F = 2$. Its symmetry group is $S_4$
(order 24), with alternating subgroup $A_4$ (order 12).

We assign gauge content to each simplex type by its boundary structure:

\begin{itemize}
\item \textit{Faces} ($p = 2$): each face has $|\partial\sigma_2|= 3$
boundary edges. The three-dimensional representation corresponds to
$\mathrm{SU}(3)$.
\item \textit{Edges} ($p = 1$): each edge has $|\partial\sigma_1|= 2$
boundary vertices. The two-dimensional representation corresponds to
$\mathrm{SU}(2)$.
\item \textit{Vertices} ($p = 0$): each vertex has one self-boundary.
The one-dimensional assignment corresponds to $\mathrm{U}(1)$.
\end{itemize}

This yields the Standard Model gauge group
$\mathrm{SU}(3)\times\mathrm{SU}(2)\times\mathrm{U}(1)$ directly
from the simplicial hierarchy of $K_4$.


% ============================================================
\section{Coupling ratios and normalization}
\label{sec:coupling}
% ============================================================

\subsection{Ratios from boundary counting}

The gauge coupling $g_p^2$ on a $p$-simplex is proportional to the
number of boundary $(p-1)$-simplices:
\begin{equation}
g_p^2 \propto |\partial\sigma_p| = p + 1 \,.
\label{eq:ratios}
\end{equation}
This determines the ratios
\begin{equation}
g_3^2 : g_2^2 : g_1^2 = 3 : 2 : 1
\label{eq:321}
\end{equation}
with no free parameters. With standard $\mathrm{SU}(5)$
normalization $g_1^2 = (5/3)\,{g'}^2$,\cite{georgi1974} the weak
mixing angle follows immediately:
\begin{equation}
\sin^2\theta_W = \frac{g'^2}{g'^2 + g_2^2}
= \frac{3/5}{3/5 + 2} = \frac{3}{13} \,.
\label{eq:sin2}
\end{equation}
This uses only the coupling ratios and is independent of the overall
normalization.

\subsection{Normalization from the Regge deficit angle}

On the $K_4$ two-skeleton---the triangulated boundary of the
tetrahedron---the $S_4$ symmetry forces equilateral faces with
interior angle $\pi/3$ at each vertex. Each vertex belongs to
$d = V - 1 = 3$ faces. The Regge deficit angle at each vertex is
\begin{equation}
\epsilon = 2\pi - d\,\frac{\pi}{3} = 2\pi - \pi = \pi \,.
\label{eq:deficit}
\end{equation}
This is exact and unique to $d = 3$: the equation
$2\pi - d\pi/3 = \pi$ has the sole integer solution $d = 3$
($V = 4$). The total curvature is $V\epsilon = 4\pi = 2\pi\chi$,
consistent with the Gauss-Bonnet theorem.

The deficit angle $\pi$ sets the coupling normalization. The
coupling on a $p$-simplex equals the boundary count divided by the
curvature quantum:
\begin{equation}
\boxed{g_p^2 = \frac{p + 1}{\pi}} \,.
\label{eq:gp}
\end{equation}


% ============================================================
\section{The strong coupling constant}
\label{sec:alphas}
% ============================================================

\subsection{Distinction coupling}

From Eq.~(\ref{eq:gp}) with $p = 2$: $g_3^2 = 3/\pi$. The
distinction coupling is
\begin{equation}
\alpha_{\mathrm{dist}} = \frac{g_3^2}{4\pi} = \frac{3}{4\pi^2}\,.
\label{eq:alphadist}
\end{equation}
Its inverse admits a geometric decomposition. The Hodge Laplacian on
$K_4$ has nonzero eigenvalues all equal to 4 (a $K_4$-specific
maximal degeneracy), with 12 nonzero modes---exactly $|A_4|$. Then
\begin{equation}
\alpha_{\mathrm{dist}}^{-1} = \frac{(2\pi\chi)^2}{|A_4|}
= \frac{16\pi^2}{12} = \frac{4\pi^2}{3} \,,
\label{eq:invdist}
\end{equation}
expressing the distinction coupling as the squared total
Gauss-Bonnet curvature divided by the number of propagating Hodge
modes.

\subsection{Support-completeness theorem}

Define the distinction number $\Delta = \binom{V}{2} + 1 = E + 1 =
7$ and the seed threshold $S_0 = (E\Delta + 1)/\Delta = 43/7$.
The product $S_0 \cdot \Delta = E^2 + E + 1 = 43$ decomposes
uniquely into three support classes:

\begin{enumerate}
\item $E^2 = 36$: edge-edge support weights. Of these, 30 arise
from direct face-sharing connections (each non-opposite edge pair
shares a face) and 6 from indirect two-step connections through the
three opposite edge pairs.
\item $E = 6$: identity-cross weights. Each edge participates in
the cell's self-identity.
\item $1$: identity-self weight. The cell distinguishes itself as a
whole---the irreducible quantum of self-distinction.
\end{enumerate}

This decomposition is exhaustive: every edge-scaffold pair is
accounted for through at most two face-mediated steps plus identity
contributions.

The spectral identity $\langle |E_{\mathrm{face}}| \rangle = E/V =
3/2$, proved from the face Hamiltonian on the $2^{2V} = 256$-state
carrier Hilbert space, provides the leading term. This identity is
specific to $K_4$: it fails for all $K_V$ with $V \geq 5$, as
verified by explicit computation.

\subsection{Result}

The strong coupling at $m_Z$ is
\begin{equation}
\boxed{\alpha_s(m_Z) = \alpha_{\mathrm{dist}} \cdot \frac{S_0}{V}
= \frac{3 \times 43}{4 \times 28 \times \pi^2}
= \frac{129}{112\pi^2}
\approx 0.1167} \,,
\label{eq:alphas}
\end{equation}
compared to the measured value $0.1180 \pm 0.0009$\cite{pdg2024}
(1.1\% deviation, within $1.4\sigma$). No renormalization group
running or matching scale enters this result.


% ============================================================
\section{Electroweak predictions}
\label{sec:ew}
% ============================================================

\subsection{Weak mixing angle}

Equation~(\ref{eq:sin2}) gives
\begin{equation}
\boxed{\sin^2\theta_W(m_Z) = \frac{3}{13} = 0.23077} \,,
\label{eq:weinberg}
\end{equation}
compared to $0.23121 \pm 0.00004$\cite{pdg2024} (0.19\% deviation).
This is the most precise prediction of the framework, requiring
only the coupling ratios (\ref{eq:321}) and standard GUT
normalization---zero free parameters.

\subsection{Fine structure constant}

The $\mathrm{SU}(2)$ coupling at $m_Z$ receives an edge support
factor $2/d = 2/3$, representing the ratio of boundary vertices per
edge to edges per vertex:
$\alpha_2(m_Z) = [g_2^2/(4\pi)] \times (2/d) = 1/(3\pi^2)$.
Combined with Eq.~(\ref{eq:weinberg}):
\begin{equation}
\boxed{\alpha_{\mathrm{EM}}^{-1}(m_Z)
= \frac{1}{\alpha_2 \sin^2\theta_W}
= 13\pi^2 \approx 128.3} \,,
\label{eq:alphaem}
\end{equation}
compared to $127.951 \pm 0.009$\cite{pdg2024} (0.28\% deviation).
The integer $13 = V + E + F - 1$ counts the proper boundary simplices
of $K_4$.

Note that $\alpha_{\mathrm{EM}}$ is not independent: it follows
algebraically from $\alpha_2$ and $\sin^2\theta_W$. The framework
therefore makes two independent gauge-sector predictions ($\alpha_s$
and $\sin^2\theta_W$) plus one consistency check
($\alpha_{\mathrm{EM}}$), all of which agree with experiment.


% ============================================================
\section{Discussion}
\label{sec:discussion}
% ============================================================

Table~\ref{tab:results} summarizes the predictions.

\begin{table}[h]
\caption{Gauge coupling predictions from $K_4$ combinatorics.
No free parameters.\label{tab:results}}
\begin{ruledtabular}
\begin{tabular}{llll}
Observable & Formula & Prediction & Measured \\
\hline
$\alpha_s(m_Z)$ & $129/(112\pi^2)$ & 0.1167 & 0.1180(9) \\
$\sin^2\theta_W(m_Z)$ & $3/13$ & 0.2308 & 0.23121(4) \\
$\alpha_{\mathrm{EM}}^{-1}(m_Z)$ & $13\pi^2$ & 128.3 & 127.951(9) \\
\end{tabular}
\end{ruledtabular}
\end{table}

Several structural features of the derivation merit comment.

\textit{Uniqueness of $K_4$.}---The vertex deficit angle equals
$\pi$ only for $d = 3$ (equivalently $V = 4$). The spectral
identity $\langle |E_{\mathrm{face}}| \rangle = E/V$ holds only
for $K_3$ (trivially) and $K_4$. The evaluation saturation theorem
---that the support weight $(E^2 + E + 1)/(E + 1)$ exactly equals
the seed threshold $(E\Delta + 1)/\Delta$ for all $K_V$---shows that
the identity self-support ($+1$) is the precise quantum needed for
any complete graph to sustain a distinction; without it, all $K_V$
fail. These three properties single out the tetrahedron among all
regular simplicial complexes.

\textit{Matter sector coupling and cubic suppression.}---The
IFS contraction ratio $F_{\mathrm{IFS}} = 1/(3\pi)$ carries a sign
that distinguishes two sectors. The geometry sector ($\chi = +2$,
forward IFS) has $F_{\mathrm{IFS}} > 0$; a matter sector related by
orientation reversal $\chi \to -\chi$ has $F_{\mathrm{IFS}} < 0$
with the same magnitude. The matter coupling is the $\mathrm{SU}(2)$
coupling distributed across all $|A_4| = 12$ IFS maps:
\begin{equation}
\alpha_{\mathrm{matter}} = \frac{\alpha_2}{|A_4|} = \frac{1}{36\pi^2}\,.
\label{eq:amatter}
\end{equation}
The ratio to the geometry (distinction) coupling is
\begin{equation}
\frac{\alpha_{\mathrm{matter}}}{\alpha_{\mathrm{dist}}} =
\frac{1}{d^3} = \frac{1}{27}\,,
\label{eq:suppression}
\end{equation}
one factor of $1/d$ per boundary dimension. Matter is suppressed
relative to geometry by the volume of the unit cube in $d$ spatial
dimensions---a purely combinatorial statement with no free parameters.

\textit{Three generations and the Koide formula from $A_4/V_4$.}---The
alternating group $A_4$ contains the Klein four-group
$V_4 = \{e,\,(12)(34),\,(13)(24),\,(14)(23)\}$ as a normal subgroup.
The three double transpositions in $V_4$ are exactly the symmetries
that preserve the pairing structure of $K_4$---the three ways to
partition four vertices into two pairs, corresponding to the three
opposite edge pairs of the tetrahedron. The quotient
\begin{equation}
A_4 / V_4 \cong Z_3
\label{eq:quotient}
\end{equation}
is cyclic of order 3. Its three cosets $C_0 = V_4$, $C_1$
(3-cycles), and $C_2$ (inverse 3-cycles) are related by the $Z_3$
generator (a $2\pi/3$ rotation around the tetrahedral axis). We
identify
\begin{equation}
C_0 \leftrightarrow e\,,\quad C_1 \leftrightarrow \mu\,,\quad
C_2 \leftrightarrow \tau\,.
\label{eq:genid}
\end{equation}

The characters of $Z_3$ are $\chi_k(C^n) = \omega^{kn}$ where
$\omega = e^{2\pi i/3}$. The Koide
parametrization~\cite{koide1983}
\begin{equation}
\sqrt{m_l} = M_K\!\left(1 + \sqrt{2}\cos\!\left(\theta +
\frac{2\pi l}{3}\right)\right), \quad l = 0, 1, 2
\label{eq:koide_param}
\end{equation}
rewrites as
\begin{equation}
\sqrt{m_l} = M_K\,\mathrm{Re}\!\left[1 + \sqrt{2}\,
e^{i\theta}\,\omega^l\right]\,,
\label{eq:koide_z3}
\end{equation}
which is the $Z_3$ character expansion of the lepton mass spectrum:
the constant term is the $\rho_0$ (trivial) contribution; the
oscillating term is the $\rho_1$ contribution with phase $e^{i\theta}$.
The three complex irreps $\rho_0$, $\rho_1$, $\rho_2$ of $Z_3$
(with $\rho_2 = \rho_1^*$ absorbed into the real part) are exactly
the three representations generated by the $A_4/V_4$ coset structure.

The Koide ratio $Q = (\sum m_l)/(\sum\sqrt{m_l})^2 = 2/3 = A_{\max}$
follows from the odd-dimensional frustration of $K_4$ (previously
derived\cite{mayes2026}). The generation count $N_{\mathrm{gen}} = 3$
follows from $|A_4/V_4| = |Z_3| = 3$. The $2\pi/3$ equal spacing of
the Koide parametrization follows from the uniform $Z_3$ characters.
All three are consequences of the group structure $A_4/V_4 \cong Z_3$
with no free parameters.

The phase $\theta$ in Eq.~(\ref{eq:koide_param}) encodes the
embedding of the mass scale into the $Z_3$ frame and requires one
mass input (equivalently, $\nu_0 = m_e/S_0$) to fix. It is not
predicted by pure group theory; it is the unique residual degree of
freedom of the lepton sector.

\textit{$S_4$ representation structure.}---The $2^{2V} = 256$-state
carrier Hilbert space decomposes under $S_4$ as
$256 = (V+1)(\Delta + d^3 + d^2 + 2^d) + 1$,
with $\mathrm{Tr}(H^2|_{\mathrm{trivial}}) = 2^5\Delta = 224$.
Every irreducible representation dimension except the sign
representation is divisible by $V + 1 = 5$, reflecting a deep
factorization tied to $K_4$ combinatorics.

\textit{Falsifiability.}---The predictions are sharp and testable.
A determination of $\alpha_s(m_Z)$ deviating from $129/(112\pi^2)$
by more than 3\%, or of $\sin^2\theta_W(m_Z)$ deviating from $3/13$
beyond $5\sigma$ of current precision, would falsify this framework.
Future lattice QCD determinations of $\alpha_s$ and precision
electroweak measurements at $e^+e^-$ colliders provide the
decisive tests.


\begin{acknowledgments}
The authors thank the broader Unified Dynamics research
collaboration for sustained discussion and cross-validation.
Computational verification was performed with AI-assisted tools.
The theoretical framework and all derivations are due to the human authors.
Supporting materials are archived at \mbox{doi:10.5281/zenodo.19582493}.
\end{acknowledgments}


\begin{thebibliography}{99}

\bibitem{pdg2024}
R.~L.~Workman \textit{et al.} (Particle Data Group),
Prog. Theor. Exp. Phys. \textbf{2022}, 083C01 (2022), and 2024
update.

\bibitem{georgi1974}
H.~Georgi and S.~L.~Glashow,
Phys.\ Rev.\ Lett.\ \textbf{32}, 438 (1974).

\bibitem{pati1974}
J.~C.~Pati and A.~Salam,
Phys.\ Rev.\ D \textbf{10}, 275 (1974).

\bibitem{polchinski1998}
J.~Polchinski,
\textit{String Theory} (Cambridge University Press, 1998).

\bibitem{koide1983}
Y.~Koide,
Phys.\ Lett.\ B \textbf{120}, 161 (1983).

\bibitem{mayes2026}
B.~W.~Mayes \textit{et al.},
Zenodo preprint, doi:10.5281/zenodo.19582493 (2026).

\end{thebibliography}

\end{document}
